\section{Wasserkraft – Hydraulische und energetische Grundlagen}

% ------------------------------------------------------------
\subsection{Kraftwerkstypen und Kennzahlen}

Typen:
\begin{itemize}
\item Laufkraftwerk
\item Speicherkraftwerk
\item Pumpspeicherkraftwerk
\item Umwälzkraftwerk
\end{itemize}

Volllaststunden:
\[
t_{VL} = \frac{E_{\text{Speicher}}}{P_{\text{Turbine}}}
\]

Leerungsfaktor:
\[
\text{LF} = \frac{\text{Zuflussmenge}}{\text{Speichervolumen}}
\]

% ------------------------------------------------------------
\subsection{Wasserdargebot und Abflusscharakteristik}

Abflussganglinie: $Q(t)$

Abflussdauerkurve: sortierte Jahresabflüsse

Definition:
\[
Q_x = \text{Abfluss, der an } x \text{ Tagen überschritten wird}
\]

Nutzwassermenge:
\[
Q = Q_B - Q_{\text{rest}}
\]

Restwasser gemäss Gesetz:
\[
Q_{\text{rest}} = f(Q_{347})
\]

% ------------------------------------------------------------
\subsection{Bernoulli-Gleichung}

Energieerhaltung:
\[
\frac{v^2}{2} + g z + \frac{p}{\rho} = \text{const}
\]

Höhengleichung:
\[
H = \frac{v^2}{2g} + z + \frac{p}{\rho g}
\]

Geschwindigkeit:
\[
v = \frac{Q}{A}
\]

% ------------------------------------------------------------
\subsection{Hydraulische Verluste}

Örtliche Verluste:
\[
h_v = \zeta \frac{v^2}{2g}
\]

Reibungsverluste (Strickler):
\[
H_{vr} =
\frac{v^2 \cdot L}
{K^2 \cdot R_h^{4/3}}
\]

Hydraulischer Radius:
\[
R_h = \frac{A}{P}
\]

Rechteck:
\[
R_h = \frac{b h}{b + 2h}
\]

Kreis:
\[
R_h = \frac{D}{4}
\]

Alternative (Moody):
\[
H_{vr} =
\lambda \frac{v^2 L}{d \cdot 2g}
\]

Reynolds-Zahl:
\[
Re = \frac{v d}{\nu}
\]

Hydraulisch rauer Bereich:
\[
\frac{1}{\sqrt{\lambda}} =
-2 \log \left( \frac{k}{3.71 d} \right)
\]

% ------------------------------------------------------------
\subsection{Nettofallhöhe}

\[
H_n = H - \sum H_{vr} - \frac{v_u^2}{2g}
\]

Quadratische Verlustabhängigkeit:
\[
\sum H_{vr} = C Q^2
\]

\[
H_n(Q) = H - C Q^2
\]

% ------------------------------------------------------------
\subsection{Leistungsberechnung}

Hydraulische Leistung:
\[
P_{hyd} = \rho Q g H_n
\]

Mechanische Leistung:
\[
P_{mech} = \eta_t P_{hyd}
\]

Elektrische Leistung:
\[
P_{el} = \eta_g \eta_t \rho Q g H_n
\]

% ------------------------------------------------------------
\subsection{Energieproduktion}

\[
E = \int P_{el} dt
\]

Diskret:
\[
E \approx \sum P_{el,i} \Delta t
\]

Wesentliche Abhängigkeiten:
\begin{itemize}
\item Verluste $\propto Q^2$
\item Leistung $\propto Q$
\item Leistung $\propto H_n$
\item Energie $\propto \int Q H_n dt$
\end{itemize}




